\documentclass[11pt,a4paper]{article}

\usepackage[utf8]{inputenc}
\usepackage[T1]{fontenc}
\usepackage[brazil]{babel}
\usepackage[margin=2.5cm]{geometry}
\usepackage{listings}
\usepackage{xcolor}
\usepackage{booktabs}
\usepackage{graphicx}
\usepackage{hyperref}
\usepackage{enumitem}
\usepackage{fancyhdr}

\hypersetup{colorlinks=true, linkcolor=blue, urlcolor=blue}

% Estilo para listagens de codigo C
\definecolor{codegray}{gray}{0.95}
\definecolor{commentgreen}{rgb}{0.0,0.5,0.0}
\definecolor{keywordblue}{rgb}{0.0,0.0,0.7}
\definecolor{stringred}{rgb}{0.6,0.0,0.0}

\lstdefinestyle{cstyle}{
    backgroundcolor=\color{codegray},
    commentstyle=\color{commentgreen},
    keywordstyle=\color{keywordblue}\bfseries,
    stringstyle=\color{stringred},
    basicstyle=\ttfamily\footnotesize,
    breaklines=true,
    captionpos=b,
    numbers=left,
    numberstyle=\tiny\color{gray},
    frame=single,
    showstringspaces=false,
    tabsize=4,
    language=C
}

\lstdefinestyle{bashstyle}{
    backgroundcolor=\color{codegray},
    basicstyle=\ttfamily\footnotesize,
    breaklines=true,
    frame=single,
    showstringspaces=false
}

\pagestyle{fancy}
\fancyhf{}
\lhead{Sistemas Operacionais}
\rhead{Prática II: Gerência de Memória}
\cfoot{\thepage}

\title{\textbf{Prática II: Gerência de Memória} \\ \large Roteiro de Laboratório: Hardware e Uso de Memória}
\author{Curso de Engenharia de Telecomunicações \\ IFSC Campus São José \\ Sistemas Operacionais (Prof. Rogério Pereira Junior)}
\date{\today}

\begin{document}

\maketitle

\section*{Introdução}
Neste laboratório investiguei o subsistema de memória do Linux. Os tópicos cobertos foram o tamanho
de página e os limites do sistema, a alocação dinâmica com \emph{Lazy Allocation}, a diferença entre
memória virtual (VSZ) e física (RSS), a inspeção da tabela de páginas com o \texttt{pmap}, as
estatísticas globais de paginação e de \emph{swap}, o mapeamento de arquivos com \texttt{mmap} e os
segmentos de memória de um processo. Todos os valores foram coletados em um ambiente Linux real
(kernel sobre WSL2, com página de 4~KB).

\tableofcontents
\vspace{1em}

% =====================================================================
\section{Roteiro Prático II: Hardware de Memória}

% ---------------------------------------------------------------------
\subsection{Tamanho da Página e Limites do Sistema}

Comandos executados:
\begin{lstlisting}[style=bashstyle]
$ getconf PAGESIZE
4096
$ sysctl kernel.shmall
kernel.shmall = 18446744073692774399
\end{lstlisting}

\paragraph{Pergunta 1: tamanho de página.}
O valor retornado foi 4096 bytes, que convertido equivale a:
\[
\frac{4096 \text{ bytes}}{1024} = 4 \text{ KB}.
\]
Esse número é o tamanho do bloco fixo (página) em que o subsistema de memória do Linux divide tanto
a memória virtual de cada processo quanto a memória física (a RAM, dividida em quadros, ou
\emph{frames}). Toda alocação e todo mapeamento de endereços virtuais para físicos é feito em
múltiplos de 4~KB, que é também a granularidade mínima de uma falta de página (\emph{page fault}).
A página é a unidade básica de gerência da MMU e da tabela de páginas.

\paragraph{Pergunta 2: limite de memória compartilhada.}
Multiplicando \texttt{kernel.shmall} pelo tamanho da página:
\[
18.446.744.073.692.774.399 \times 4096 \text{ bytes} \approx 7{,}03 \times 10^{22} \text{ bytes}.
\]
Em GB isso corresponde a aproximadamente $7 \times 10^{13}$~GB. Na prática, esse é o valor padrão do
kernel (próximo do limite teórico de um inteiro de 64 bits), o que significa que o \texttt{shmall}
não impõe limite prático para alocação de memória compartilhada neste sistema. O fator limitante
real passa a ser a RAM física e o espaço de \emph{swap} disponíveis. Em sistemas onde o administrador
define um \texttt{shmall} menor, esse produto fornece o teto real de memória compartilhada.

% ---------------------------------------------------------------------
\subsection{Alocação Dinâmica e Páginas Virtuais vs. Físicas}

O programa \texttt{paginacao.c} aloca 16~MB com \texttt{malloc} e só escreve neles em um segundo
momento, o que permite observar o efeito da \emph{Lazy Allocation}.

\lstinputlisting[style=cstyle,caption={paginacao.c}]{paginacao.c}

% ---------------------------------------------------------------------
\subsection{Compilação e Monitoramento (\texttt{ps} e \texttt{/proc})}

\begin{lstlisting}[style=bashstyle]
$ gcc paginacao.c -o paginacao
$ ./paginacao            # Terminal 1
# Terminal 2:
$ ps -o pid,vsz,rss,comm -p [PID]
$ cat /proc/[PID]/status | grep -E "VmSize|VmRSS"
\end{lstlisting}

Valores coletados em cada etapa, com o PID monitorado durante a execução real:

\begin{table}[h]
\centering
\begin{tabular}{lrr}
\toprule
Momento do Programa & VSZ / VmSize (KB) & RSS / VmRSS (KB) \\
\midrule
Passo A (Início)        & 2.684  & 1.408 \\
Passo B (malloc feito)  & 19.072 & 1.408 \\
Passo C (Escrita feita) & 19.072 & 17.792 \\
\bottomrule
\end{tabular}
\caption{Memória virtual (VSZ) e física residente (RSS) por etapa.}
\end{table}

\paragraph{Pergunta: o que ocorreu no Passo B (malloc)?}
A memória virtual (VSZ) cresceu de 2.684~KB para 19.072~KB, um aumento de cerca de 16.384~KB
($=16$~MB), que corresponde à região reservada pelo \texttt{malloc}. Já a memória física (RSS)
permaneceu inalterada em 1.408~KB. O motivo é a \emph{Lazy Allocation} (alocação preguiçosa): o
\texttt{malloc} apenas reserva um intervalo de endereços no espaço virtual do processo, ajustando a
tabela de páginas com entradas marcadas, mas o kernel não associa nenhum quadro físico de RAM
enquanto aquelas páginas não forem efetivamente acessadas. Como ainda não houve escrita, nenhum
\emph{page fault} ocorreu e a RAM não foi consumida.

\paragraph{Pergunta: o que ocorreu no Passo C (escrita) e por que pular de 4096 em 4096?}
Depois do laço \texttt{for}, o RSS saltou de 1.408~KB para 17.792~KB (cerca de $+16$~MB), mostrando
que agora os 16~MB virtuais passaram a ocupar quadros físicos reais na RAM. Isso acontece porque a
primeira escrita em cada página dispara uma falta de página (\emph{page fault}); o kernel então
aloca um quadro físico livre, atualiza a tabela de páginas mapeando o endereço virtual ao físico e
só então a escrita se conclui.

O código pula de 4096 em 4096 bytes porque esse é o tamanho da página. Basta tocar um único byte de
cada página para forçar a alocação do quadro inteiro de 4~KB; escrever em todos os bytes seria
desperdício de tempo, pois não traz nenhuma página física adicional. Com o passo de 4096 garantimos
que todas as páginas dos 16~MB sejam materializadas na RAM, percorrendo o número mínimo de posições
($16\,\text{MB}/4\,\text{KB}=4096$ páginas).

% ---------------------------------------------------------------------
\subsection{Inspecionando a Tabela de Páginas com \texttt{pmap}}

Com o programa parado no Passo C, executei \texttt{pmap -x [PID]}. A linha correspondente à alocação
de 16~MB:

\begin{lstlisting}[style=bashstyle]
Address           Kbytes     RSS   Dirty Mode  Mapping
00007fcf4a1ff000   16388   16384   16384 rw---   [ anon ]
total kB           19072   17936   16488
\end{lstlisting}

\begin{table}[h]
\centering
\begin{tabular}{lr}
\toprule
Coluna & Valor (KB) \\
\midrule
Kbytes (tamanho virtual da região) & 16.388 \\
RSS (páginas residentes em RAM)    & 16.384 \\
Dirty (páginas sujas)              & 16.384 \\
\bottomrule
\end{tabular}
\caption{Região anônima de 16~MB observada no \texttt{pmap -x}.}
\end{table}

\paragraph{Pergunta: o que a coluna Dirty indica?}
A coluna Dirty (página suja) indica a quantidade de memória cujas páginas foram modificadas
(escritas) após o mapeamento e cujo conteúdo na RAM difere da fonte em disco, ou que nunca teve
respaldo em disco, como é o caso da memória anônima. Para o Sistema Operacional, uma página suja é
uma página que não pode ser simplesmente descartada: antes de reaproveitar seu quadro físico, o
kernel precisa primeiro gravá-la, seja no \emph{swap} (no caso de memória anônima) seja de volta no
arquivo (no caso de páginas mapeadas em arquivo). Uma página limpa, idêntica ao disco, pode ser
descartada na hora, pois é possível relê-la depois. Como escrevemos \texttt{'X'} em todas as 4096
páginas, todas as 16.384~KB aparecem como sujas.

% ---------------------------------------------------------------------
\subsection{Páginas Globais e Atividade de \emph{Swap} (\texttt{/proc/meminfo} e \texttt{vmstat})}

\begin{lstlisting}[style=bashstyle]
$ cat /proc/meminfo | grep -E "Cached|Buffers"
Buffers:           87204 kB
Cached:          1933848 kB
SwapCached:            0 kB

$ vmstat 2
procs -----------memory---------- ---swap-- -----io---- -system-- ------cpu-----
 r  b   swpd    free   buff   cache   si  so   bi   bo   in   cs us sy id wa st
 0  0      0 12326784  87344 2028772    0   0 1465 1849 2146    2  1  1 98  0  0
 0  0      0 12326052  87344 2028780    0   0    0    0  726 1251  0  0 100  0  0
\end{lstlisting}

\paragraph{Pergunta: o que indicam as colunas \texttt{si} (swap in) e \texttt{so} (swap out)?}
Sob o cabeçalho \texttt{swap} do \texttt{vmstat}:
\begin{itemize}[nosep]
    \item \texttt{si} (swap in): quantidade de memória, em KB/s, lida do disco de \emph{swap} de
    volta para a RAM. Indica que páginas que haviam sido enviadas ao \emph{swap} estão sendo trazidas
    de novo porque algum processo precisou delas.
    \item \texttt{so} (swap out): quantidade de memória, em KB/s, escrita da RAM para o disco de
    \emph{swap}. Ocorre quando há pressão de memória e o kernel precisa liberar quadros físicos,
    movendo páginas pouco usadas para o \emph{swap}.
\end{itemize}
Valores altos e contínuos nessas colunas indicam \emph{thrashing}, ou seja, um sistema com falta de
RAM e desempenho degradado. Nos testes, ambos ficaram em \texttt{0}, já que o sistema tinha RAM livre
de sobra (cerca de 11~GiB).

\paragraph{Pergunta: o que são \emph{HugePages}?}
\begin{lstlisting}[style=bashstyle]
$ cat /proc/meminfo | grep -i huge
AnonHugePages:         0 kB
HugePages_Total:       0
HugePages_Free:        0
Hugepagesize:       2048 kB
\end{lstlisting}
HugePages são páginas de memória de tamanho muito maior do que a página padrão de 4~KB, tipicamente
2~MB (visível em \texttt{Hugepagesize: 2048 kB}) ou 1~GB em arquiteturas x86-64. O objetivo é reduzir
a sobrecarga de gerência de memória: com páginas maiores, são necessárias muito menos entradas na
tabela de páginas para cobrir a mesma quantidade de RAM, o que diminui o número de \emph{misses} na
TLB (\emph{Translation Lookaside Buffer}, a cache de traduções da MMU) e acelera o acesso à memória.
São úteis para aplicações que manipulam grandes volumes contíguos de memória, como bancos de dados
(Oracle, PostgreSQL) e máquinas virtuais. Os campos \texttt{HugePages\_Total} e
\texttt{HugePages\_Free} mostram quantas estão pré-reservadas (zero neste sistema), e
\texttt{AnonHugePages} reflete o uso de \emph{Transparent Huge Pages} (THP), em que o kernel promove
páginas automaticamente.

% ---------------------------------------------------------------------
\subsection{Mapeamento de Arquivos em Páginas (\texttt{mmap})}

\lstinputlisting[style=cstyle,caption={mapeado.c}]{mapeado.c}

\begin{lstlisting}[style=bashstyle]
$ echo "Arquitetura de Computadores e Sistemas Operacionais" > texto.txt
$ gcc mapeado.c -o mapeado
$ ./mapeado
Conteudo lido da pagina de memoria: Arquitetura de Computadores e Sistemas Operacionais
$ cat texto.txt
Orquitetura de Computadores e Sistemas Operacionais
\end{lstlisting}

\paragraph{Pergunta: o que mudou em \texttt{texto.txt} e por quê?}
O primeiro caractere mudou de \texttt{'A'} para \texttt{'O'}: o arquivo passou de
"\texttt{Arquitetura...}" para "\texttt{Orquitetura...}". A alteração veio da única linha
\texttt{arquivo\_em\_memoria[0] = 'O';}, que é uma escrita em uma posição de memória, sem nenhuma
chamada a \texttt{write} ou \texttt{fprintf}.

Isso é possível por causa do \texttt{mmap} com a \emph{flag} \texttt{MAP\_SHARED}, que mapeia o
conteúdo do arquivo em disco diretamente no espaço de endereçamento virtual do processo. As páginas
virtuais retornadas pelo \texttt{mmap} ficam associadas ao arquivo. Quando o processo escreve em
\texttt{arquivo\_em\_memoria[0]}, a MMU traduz o endereço virtual para o físico e a página
correspondente é marcada como suja. O kernel, por meio do \emph{page cache} e do \emph{writeback},
sincroniza automaticamente as páginas sujas de volta ao arquivo em disco, seja no \texttt{munmap},
seja em um \emph{flush} periódico, seja sob pressão de memória. Dessa forma, a paginação transforma
operações de E/S de arquivo em simples acessos à memória, sem precisar de chamadas explícitas de
escrita.

% ---------------------------------------------------------------------
\subsection{Monitoramento de paginação em tempo real (\texttt{sar -B})}

O comando \texttt{sar -B 1 10} (do pacote \texttt{sysstat}) coleta estatísticas de paginação a cada
1~segundo, 10 vezes. As principais colunas e seus significados:
\begin{itemize}[nosep]
    \item pgpgin/s, pgpgout/s: KB/s lidos do disco para a memória e gravados da memória para o
    disco, respectivamente.
    \item fault/s: número total de faltas de página por segundo (\emph{minor} mais \emph{major}).
    \item majflt/s: faltas de página maiores por segundo, aquelas que exigem acesso ao disco (a
    página precisa ser carregada de um arquivo ou do \emph{swap}) e que são as mais custosas.
    \item pgfree/s: páginas liberadas pelo sistema por segundo.
    \item pgscank/s, pgscand/s: páginas varridas pelo \emph{kswapd} (em segundo plano) e diretamente,
    em busca de páginas para recuperar.
    \item pgsteal/s: páginas recuperadas da \emph{page cache} por segundo.
    \item \%vmeff: eficiência da recuperação de páginas (\texttt{pgsteal/pgscan}); valores próximos
    de 100\% indicam recuperação eficiente.
\end{itemize}
Observação: neste ambiente o pacote \texttt{sysstat} não estava instalado
(\texttt{sudo apt install sysstat}), então a descrição acima cobre o significado de cada linha.

% =====================================================================
\section{Roteiro Prático II: Uso de Memória}

% ---------------------------------------------------------------------
\subsection{Segmentos de um processo (\texttt{hello.c} e \texttt{pmap})}

\lstinputlisting[style=cstyle,caption={hello.c}]{hello.c}

A variável \texttt{global\_init = 42} fica no segmento DATA (dados inicializados) e
\texttt{global\_uninit} no segmento BSS (dados não inicializados, zerados na carga). A saída do
\texttt{pmap} mostra os segmentos do processo:

\begin{lstlisting}[style=bashstyle]
$ ./hello & ; pmap <PID>
4646:   ./hello
000064557c8dc000      4K r---- hello          <- .text/rodata (codigo, RO)
000064557c8dd000      4K r-x-- hello           <- codigo executavel
000064557c8de000      4K r---- hello
000064557c8df000      4K r---- hello           <- DATA / BSS
000064557c8e0000      4K rw--- hello           <- dados graváveis
00006455a39f7000    132K rw---   [ anon ]      <- HEAP
000074ca99e00000   ...  libc.so.6              <- biblioteca C compartilhada
000074ca9a1d1000   ...  ld-linux-x86-64.so.2   <- linker dinâmico
00007ffebcd64000    136K rw---   [ stack ]      <- STACK (pilha)
 total             2684K
\end{lstlisting}

Dá para ver a divisão em segmentos: código (permissões \texttt{r-x}, somente leitura e execução),
dados (\texttt{rw-}), heap (\texttt{[ anon ]}), pilha (\texttt{[ stack ]}) e as bibliotecas
compartilhadas (\texttt{libc} e \texttt{ld-linux}). Todas as regiões são múltiplos de 4~KB, o
tamanho da página.

% ---------------------------------------------------------------------
\subsection{Alocação estática, automática e dinâmica (\texttt{alocacao.c})}

\lstinputlisting[style=cstyle,caption={alocacao.c}]{alocacao.c}

\begin{lstlisting}[style=bashstyle]
$ ./alocacao
Local = 5
Dinamica = 99
\end{lstlisting}

\noindent Os três tipos de alocação:
\begin{itemize}[nosep]
    \item Estática (\texttt{int global = 10}): vive no segmento DATA e existe durante toda a execução
    do programa.
    \item Automática (\texttt{int local = n}): vive na pilha (STACK), criada na entrada da função e
    destruída no seu retorno.
    \item Dinâmica (\texttt{malloc}): vive no heap e é controlada manualmente pelo programador, com
    \texttt{malloc} e \texttt{free}.
\end{itemize}

% ---------------------------------------------------------------------
\subsection{Recursão e a Pilha (\texttt{fatorial.c})}

\lstinputlisting[style=cstyle,caption={fatorial.c}]{fatorial.c}

\begin{lstlisting}[style=bashstyle]
$ ./fatorial
inicio: n = 4
inicio: n = 3
inicio: n = 2
inicio: n = 1
final : n = 1, parcial = 1
final : n = 2, parcial = 2
final : n = 3, parcial = 6
final : n = 4, parcial = 24
Fatorial(4) = 24
\end{lstlisting}

A saída mostra o comportamento de pilha (LIFO). Todas as chamadas "\texttt{inicio}" aparecem
primeiro, na descida da recursão ($n=4,3,2,1$): cada chamada de \texttt{fatorial} empilha um novo
\emph{stack frame} com a sua própria variável local \texttt{parcial} e o parâmetro \texttt{n}, como
pratos empilhados. Ao chegar no caso base ($n<2$), os retornos acontecem na ordem inversa
($n=1,2,3,4$), desempilhando cada quadro e calculando o produto parcial. Cada nível mantém a sua
cópia das variáveis locais, o que mostra que a pilha isola o contexto de cada chamada.

% =====================================================================
\section{Considerações Finais: a função da MMU}

A MMU (Memory Management Unit) é o componente de hardware, em geral integrado ao processador, que
traduz em tempo de execução os endereços virtuais usados pelos programas em endereços físicos reais
da RAM. Toda referência de memória feita por um processo passa pela MMU, que consulta a tabela de
páginas mantida pelo Sistema Operacional para mapear a página virtual ao quadro físico
correspondente, usando a cache TLB para acelerar as traduções recentes.

Com base nos experimentos, as funções da MMU que observei foram:
\begin{itemize}[nosep]
    \item Tradução de endereços virtuais para físicos, na granularidade de página (4~KB), de forma
    transparente para o programa.
    \item Suporte à \emph{Lazy Allocation}: ao detectar acesso a uma página virtual ainda sem quadro
    físico, a MMU gera uma falta de página (\emph{page fault}) e interrompe a CPU para que o kernel
    aloque o quadro. Foi exatamente o que vimos entre os Passos B e C.
    \item Proteção e isolamento de memória: a MMU verifica as permissões de cada página (leitura,
    escrita e execução) e impede que um processo acesse a memória de outro.
    \item Marcação de páginas como sujas ou acessadas, dando ao SO a informação necessária para
    decidir o que enviar ao \emph{swap} e o que sincronizar com o disco, como no \texttt{mmap}.
\end{itemize}

A MMU é, portanto, a peça que viabiliza a memória virtual: ela separa o espaço de endereçamento
lógico de cada processo da memória física e permite paginação, proteção, compartilhamento de memória
e uso de \emph{swap}, tudo de forma transparente para as aplicações.

\end{document}
